\section{Use an impulse to expose order and scale}
The supplied caller has seven zero history samples and an impulse code 16000 at the
first new sample. Its eight coefficients are 4096. For the first eight outputs,
exactly one tap sees the impulse, so each output is 2000. Later windows contain only
zeros. With default count sixteen, the result is eight samples of 2000 followed by
eight zeros, with no clipping.

Equal coefficients cannot reveal reversed order. For a second impulse test, use
code 16384 and oldest-first coefficients
\(1024,2048,3072,4096,5120,6144,7168,8192\). The first current sample meets the last
stored coefficient. The output sequence is therefore
\(4096,3584,3072,2560,2048,1536,1024,512\), followed by zeros. These products need
no rounding approximation. Reversing this response indicates a coefficient-order error.

An impulse test checks more than whether the output eventually becomes zero. Check
which output receives each coefficient, the scale of every nonzero value, the seven
history positions and the point where the impulse leaves the window. It does not by
itself establish behavior for negative values, large sums or rounding ties; those
require separate examples.
